\documentclass{article}
\usepackage{neurips_2025}
\usepackage[utf8]{inputenc}
\usepackage[T1]{fontenc}
\usepackage{hyperref}
\usepackage{booktabs}
\usepackage{amsfonts}
\usepackage{nicefrac}
\usepackage{microtype}
\usepackage{xcolor}
\usepackage{pgfplots}
\pgfplotsset{compat=1.18}

\title{DSAA2011 Course Project: Student Dropout Prediction and Clustering Analysis}
\author{Group Project Report \\
DSAA2011 (Spring 2026)}

\begin{document}
\maketitle

\begin{abstract}
We complete the DSAA2011 project tasks on the Student Dropout dataset (\texttt{data.csv}). The workflow covers preprocessing, low-dimensional visualization (PCA and t-SNE), clustering, supervised training/testing, and model evaluation. We implement all methods in pure Python due environment constraints and report performance using confusion-matrix-based metrics and ROC-AUC. Logistic Regression achieves the strongest generalization among tested models, while K-Means outperforms K-Medoids for unsupervised grouping quality.
\end{abstract}

\section{Introduction}
Predicting student dropout risk is an important educational analytics problem with direct intervention value. We use the Student Dropout dataset containing 4,424 students and 36 attributes (demographics, background, and academic indicators) with three labels: \texttt{Dropout}, \texttt{Graduate}, and \texttt{Enrolled}. 

For supervised learning, we define a binary target: \texttt{Dropout}=1 and \texttt{Graduate/Enrolled}=0. Our objectives are:
\begin{itemize}
    \item produce a robust preprocessing pipeline,
    \item explore structure through low-dimensional projection and clustering,
    \item train and compare at least two predictive models,
    \item evaluate with accuracy, precision, recall, F1, ROC and AUC,
    \item provide model choice and further discussion.
\end{itemize}

\section{Mandatory Tasks}
\paragraph{Implementation update.} In the refined submission, the primary technical artifact is a Jupyter notebook (\texttt{project.ipynb}) implemented with pandas/numpy/matplotlib/seaborn/scikit-learn to provide polished figures, multiclass confusion matrices, and per-class performance tables for all target categories.

\subsection{Data preprocessing}
All 36 predictors were parsed as numeric values from the semicolon-separated source file. Missing values were checked programmatically (none found in parsed features). We standardized all features with z-score normalization to avoid scale bias in distance-based and gradient-based models.

The class distribution is imbalanced: Dropout 1,421; Graduate 2,209; Enrolled 794. This motivates using precision/recall/F1 and ROC-AUC in addition to plain accuracy.

\subsection{Data visualization}
We generated 2D embeddings for visual analysis (file: \texttt{results/embedding\_points.csv}), including a deterministic PCA projection and a custom t-SNE embedding on a runtime-feasible subsample. The plots show partial class separation with overlap regions, suggesting that dropout risk is not linearly separable in all areas.

\subsection{Clustering analysis}
We compared two clustering algorithms (k=3) on a representative standardized subset:
\begin{itemize}
    \item K-Means
    \item K-Medoids
\end{itemize}

\begin{table}[h]
\centering
\caption{Clustering quality metrics (higher silhouette, lower DB is better).}
\begin{tabular}{lcc}
\toprule
Algorithm & Silhouette & Davies--Bouldin \\
\midrule
K-Means & 0.2292 & 1.5734 \\
K-Medoids & 0.1711 & 1.8603 \\
\bottomrule
\end{tabular}
\end{table}

K-Means consistently provides more compact and separated clusters. K-Medoids is more robust to outliers in principle, but here it underfits cluster structure.

\subsection{Prediction: training and testing}
We used stratified 80/20 split and compared:
\begin{itemize}
    \item Logistic Regression (L2 regularization)
    \item k-Nearest Neighbors (k=17)
\end{itemize}

\begin{table}[h]
\centering
\caption{Predictive performance summary.}
\begin{tabular}{llccccc}
\toprule
Model & Split & Acc & Prec & Rec & F1 & AUC \\
\midrule
LR & Train & 0.8720 & 0.8708 & 0.7060 & 0.7798 & 0.9155 \\
LR & Test & 0.8725 & 0.8644 & 0.7158 & 0.7831 & 0.9066 \\
LR & All & 0.8721 & 0.8695 & 0.7080 & 0.7804 & 0.9138 \\
KNN & Test & 0.7844 & 0.8092 & 0.4316 & 0.5629 & 0.8365 \\
\bottomrule
\end{tabular}
\end{table}

Test confusion matrix values:
\begin{itemize}
    \item LR: TP=204, TN=569, FP=32, FN=81
    \item KNN: TP=123, TN=572, FP=29, FN=162
\end{itemize}

Interpretation: KNN has very high precision but low recall (misses many actual dropout students). LR has better recall/F1 trade-off and stronger AUC, which is more suitable for early-risk screening.

\subsection{Evaluation and model choice}
Using confusion-matrix metrics and ROC-AUC, we choose \textbf{Logistic Regression} as the primary model for this dataset. It provides stable train-test behavior (small overfitting gap), higher recall, and better balanced performance under class imbalance. A lightweight cross-validation study also shows LR performance is stable for L2 in the range $10^{-4}$ to $10^{-2}$, while sampled kNN validation favors $k=17$.

\section{Open-ended Exploration}
We explored model/algorithm comparisons beyond one baseline:
\begin{itemize}
    \item Unsupervised: K-Means vs K-Medoids with two internal metrics.
    \item Supervised: LR vs KNN with ROC-AUC and confusion matrices.
\end{itemize}

Findings suggest that a next-step extension should include non-linear ensemble models (e.g., random forest/gradient boosting) and probability calibration to increase recall without excessive false positives.

\section{Conclusion}
This project completed all core DSAA2011 tasks on the Student Dropout dataset: preprocessing, visualization, clustering, supervised modeling, and model evaluation. Logistic Regression is selected as the best overall predictor in current experiments. K-Means yields better clustering quality than K-Medoids. The resulting pipeline is reproducible and all generated outputs are saved under \texttt{results/}.

\section*{References}
\begin{itemize}
    \item Project announcement PDF provided by course staff (DSAA2011-26sp-project\_announce-L01.pdf).
    \item Dataset: Student Dropout Dataset (\texttt{data.csv}, provided in project repository).
\end{itemize}

\section*{Credit and GenAI disclosure}
This report text and code organization were polished with assistance from an AI coding assistant. Core dataset processing, model logic, and all quantitative outputs are produced within this repository's scripts.

\end{document}
